% Artículo VII — primera redacción editorial, 10 de septiembre de 2026.
% Apertura integrada con las demostraciones de la revisión 06.
% Procedencia y condiciones: PLAN_RESIDENCIAS_DEMOSTRATIVAS.md.
% Las referencias internas remiten a demostraciones incluidas en el cuerpo.

\begin{center}\bfseries Resumen\end{center}
\begingroup
\fontsize{9.5}{10.7}\selectfont
\setlength{\parskip}{1.5pt plus .5pt minus .5pt}
\noindent
Se determina la amplitud torsional como funcional explícito del estado
material mediante la acción espinorial, la eliminación de la contorsión
y la variación métrica. Se construye un sector material en el que esa
amplitud es positiva y se conserva. Su contribución energética produce
una dilución \(a^{-6}\) y, en la realización homogénea plana con
\(\Lambda_{\mathrm{eff}}=0\), densidades positivas y \(w<1\), un umbral
de rebote único y transversal. Para polvo se obtiene una solución exacta
para todo tiempo, con curvatura finita y geodésicas causales de
Levi--Civita completas en ambas direcciones.

La construcción parte de un núcleo formal denominado Holografía Modular
Triádica, en adelante HMT. \emph{All Possible Paths}, en adelante APP,
reúne caminos aritméticos sobre \((\mathbb Z/9\mathbb Z)^2\), con
adyacencia dada por el grafo de Cayley aditivo de generadores cardinales
y realización celular toroidal, distinguiendo las operaciones y las
evaluaciones de suma y producto;
TRIT es la unidad de firma orientada \(\{-1,0,+1\}\) que determina los
regímenes locales; el \emph{Topological Phase Kernel}, en adelante TPK,
compone la dinámica de transporte y memoria. El estado enriquecido
conserva hojas, orientación, residuo, acarreo e historia durante el
retorno de fase. La Mecánica Dimensional expresa sus resultados
internos mediante magnitudes y observables físicos, preservando la
procedencia estructural de los acoplamientos.

La normalización gravitatoria se desarrolla previamente mediante la
composición longitudinal \(L=(5759/23040)\alpha^{16}L_*\) y su
métrica radial positiva. En la realización definida por esas operaciones
y por la identificación del radio gravitatorio reducido, la reciprocidad
con la sección de acción fija \(G=c^3L^2/\hbar_{\rm ret}\) en todos
los modos admisibles. La demostración conserva el archivo completo y
establece las condiciones de reducción estática y dinámica de la memoria.
La elección dimensional \(L_*=1\,\mathrm m\), las reglas longitudinal
y métrica y la identificación radial se declaran en el propio enunciado.

El transporte de rutas admite una realización de Cartan compatible con
composición, inversión y refinamiento. Sobre ella se distinguen dos
acciones materiales: la extensión minimizante de la pantalla de memoria
y la realización espinorial. Cada una conserva su propia variación
respecto de la conexión. En el sector espinorial, afín en la contorsión,
la inversión de los operadores de Holst y torsión determina la corrección
efectiva; para la extensión mínima se conserva la dependencia no lineal
de la corriente y se demuestra un criterio local mediante el Hessiano.
La corriente, el observable cuártico y su normalización explican el
valor de la amplitud, además de su papel en la evolución cosmológica.

La persistencia del rebote bajo perturbaciones se demuestra mediante
cotas locales diferenciadas, mientras la completitud causal corresponde
a la solución global explícita de polvo. Los balances residuales, la
información de horizonte y las operaciones del observador conservan las
condiciones de fuentes, acoplamientos y coordenatizaciones que utilizan. La respuesta
relacional del estado conjunto se enlaza con la propagación avanzada y
retardada mediante operadores de Green para fuentes de soporte finito,
asociados a la diferencia covariante
construida a partir del transporte de rutas y una condición de
compatibilidad en la frontera. El resultado relaciona la memoria
generada, su realización geométrica y las consecuencias dinámicas e
informacionales de esa realización.

La composición entre acción, incremento areal de Barbero--Immirzi y
horizonte determina la conversión
\(T_{\rm cos}(L_\alpha)/\Theta_{\rm clk}=\log3/(2\pi)\).
El Hamiltoniano de respuesta angular proporciona un balance finito de
energía libre mediante la entropía relativa y una primera ley a radio
fijo. Una sección termométrica relativa caracteriza, además, el
transductor positivo y el criterio para identificar otro lector sobre
esa misma sección.

\noindent\textbf{Palabras clave:} gravitación; torsión; memoria aritmética; ecuaciones de campo; longitud de Planck; transición de contracción a expansión.
\par\endgroup
